% GENERATED FILE. Edit canonical lessons, then run npm run generate:books.
% canonical-source-hash: fnv1a64:35c44e1cc4f64d03
% canonical-lessons: TA-C356-kalanturaiyatu, TA-C356-peccuvarttai, TA-C356-kurukkitu, TA-C356-camaracam, TA-C356-vakkuvatam

\chapter{Talking Things Over}
\label{ch:ta-talking-things-over}
\hlchaptermodality{356}

\begin{chapteropening}
\textbf{By the end of this chapter:} \emph{I can talk about discussions, interruptions and compromise.}

Complete the last lesson of chapter 356: I can talk about discussions, interruptions and compromise.
\end{chapteropening}

\section[\texorpdfstring{\ta{கலந்துரையாடு} (kalanturaiyāṭu)}{கலந்துரையாடு (kalanturaiyāṭu)}]{\ta{கலந்துரையாடு} (kalanturaiyāṭu) — to talk over; to discuss together}

\label{lesson:TA-C356-kalanturaiyatu}



\begin{quote}
Before the new one: say the Tamil for to consult, then the Tamil for to propose.
\end{quote}

\subsection*{You'll want to know: \ta{கலந்துரையாடு}}
\textbf{\ta{கலந்துரையாடு}} — \emph{kalanturaiyāṭu} — \textquotedblleft{}to talk over; to discuss together\textquotedblright{}.

\begin{quote}
\ta{மாணவர்கள்} \ta{இந்தத்} \ta{தலைப்பைப்} \ta{பற்றிக்} \ta{கலந்துரையாடினார்கள்}.
\end{quote}

\begin{quote}
\emph{māṇavarkaḷ intat talaippaip paṟṟik kalanturaiyāṭiṉārkaḷ.}
\end{quote}

\begin{quote}
The students talked over this topic.
\end{quote}

Past (I …): \ta{கலந்துரையாடினேன்}, \emph{kalanturaiyāṭiṉēṉ}.

The topic takes -\ta{ஐப்} \ta{பற்றி}; the others -\ta{உடன்}.

\begin{cousinweb}
\textbf{\ta{கலந்துரையாடு}} (\emph{kalanturaiyāṭu}) is native: \emph{kalantu}, \textquotedblleft{}having mixed, joined\textquotedblright{}, from \emph{kala} (DEDR 1299), \emph{urai}, \textquotedblleft{}word\textquotedblright{} (DEDR 648), and \emph{āṭu}, \textquotedblleft{}to move, play\textquotedblright{} (DEDR 347): to exchange words together.
\end{cousinweb}

\begin{grammarlens}[title={What you've built}]
One more word for giving reasons and opinions.
\end{grammarlens}

\subsection*{Your turn}
\begin{itemize}
\raggedright
  \item \emph{Say it:} \emph{kalanturaiyāṭu}
  \item \emph{Say it:} \emph{kalanturaiyāṭu}, once more
  \item \emph{Say it:} \emph{māṇavarkaḷ intat talaippaip paṟṟik kalanturaiyāṭiṉārkaḷ.} — The students talked over this topic.
  \item \emph{Recall:} say the Tamil for to consult, then the Tamil for to propose, then say \emph{kalanturaiyāṭu} again
\end{itemize}

\begin{tcolorbox}[breakable,colback=teal!4,colframe=teal!35!black,title={Before you move on}]
What does \ta{கலந்துரையாடு} mean? (\textquotedblleft{}To talk over; to discuss together\textquotedblright{}.) Say it once more.
\end{tcolorbox}
\section[\texorpdfstring{\ta{பேச்சுவார்த்தை} (pēccuvārttai)}{பேச்சுவார்த்தை (pēccuvārttai)}]{\ta{பேச்சுவார்த்தை} (pēccuvārttai) — negotiations; talks}

\label{lesson:TA-C356-peccuvarttai}



\begin{quote}
Before the new one: say the Tamil for to propose, then the Tamil for to talk over.
\end{quote}

\subsection*{You'll want to know: \ta{பேச்சுவார்த்தை}}
\textbf{\ta{பேச்சுவார்த்தை}} — \emph{pēccuvārttai} — \textquotedblleft{}negotiations; talks\textquotedblright{}.

\begin{quote}
\ta{இரண்டு} \ta{நாடுகளுக்கும்} \ta{இடையே} \ta{பேச்சுவார்த்தை} \ta{நடக்கிறது}.
\end{quote}

\begin{quote}
\emph{iraṇṭu nāṭukaḷukkum iṭaiyē pēccuvārttai naṭakkiṟatu.}
\end{quote}

\begin{quote}
Talks are going on between the two countries.
\end{quote}

\ta{பேச்சுவார்த்தை} \ta{நடக்கிறது}, talks are under way.

\begin{cousinweb}
\textbf{\ta{பேச்சுவார்த்தை}} (\emph{pēccuvārttai}) joins native \emph{pēccu}, \textquotedblleft{}speech\textquotedblright{}, from \emph{pēcu}, \textquotedblleft{}to talk\textquotedblright{} (DEDR 4430), and \emph{vārttai}, \textquotedblleft{}word\textquotedblright{}, borrowed from Sanskrit \emph{vārttā}, \textquotedblleft{}talk, tidings\textquotedblright{} (Monier-Williams).
\end{cousinweb}

\begin{grammarlens}[title={What you've built}]
One more word for giving reasons and opinions.
\end{grammarlens}

\subsection*{Your turn}
\begin{itemize}
\raggedright
  \item \emph{Say it:} \emph{pēccuvārttai}
  \item \emph{Say it:} \emph{pēccuvārttai}, once more
  \item \emph{Say it:} \emph{iraṇṭu nāṭukaḷukkum iṭaiyē pēccuvārttai naṭakkiṟatu.} — Talks are going on between the two countries.
  \item \emph{Recall:} say the Tamil for to propose, then the Tamil for to talk over, then say \emph{pēccuvārttai} again
\end{itemize}

\begin{tcolorbox}[breakable,colback=teal!4,colframe=teal!35!black,title={Before you move on}]
What does \ta{பேச்சுவார்த்தை} mean? (\textquotedblleft{}Negotiations; talks\textquotedblright{}.) Say it once more.
\end{tcolorbox}
\section[\texorpdfstring{\ta{குறுக்கிடு} (kuṟukkiṭu)}{குறுக்கிடு (kuṟukkiṭu)}]{\ta{குறுக்கிடு} (kuṟukkiṭu) — to interrupt; to intervene}

\label{lesson:TA-C356-kurukkitu}



\begin{quote}
Before the new one: say the Tamil for to talk over, then the Tamil for negotiations.
\end{quote}

\subsection*{You'll want to know: \ta{குறுக்கிடு}}
\textbf{\ta{குறுக்கிடு}} — \emph{kuṟukkiṭu} — \textquotedblleft{}to interrupt; to intervene\textquotedblright{}.

\begin{quote}
\ta{நான்} \ta{பேசும்போது} \ta{அவன்} \ta{குறுக்கிட்டான்}.
\end{quote}

\begin{quote}
\emph{nāṉ pēcumpōtu avaṉ kuṟukkiṭṭāṉ.}
\end{quote}

\begin{quote}
He interrupted while I was speaking.
\end{quote}

Past (I …): \ta{குறுக்கிட்டேன்}, \emph{kuṟukkiṭṭēṉ}.

Polite: \ta{குறுக்கிடுவதற்கு} \ta{மன்னிக்கவும்}, sorry to interrupt.

\begin{cousinweb}
\textbf{\ta{குறுக்கிடு}} (\emph{kuṟukkiṭu}) is native: \emph{kuṟukku}, \textquotedblleft{}shortness, crossing\textquotedblright{}, from \emph{kuṟu}, \textquotedblleft{}short\textquotedblright{} (DEDR 1851), and \emph{iṭu}, \textquotedblleft{}to put\textquotedblright{} (DEDR 442): to put oneself across someone's path.
\end{cousinweb}

\begin{grammarlens}[title={What you've built}]
One more word for giving reasons and opinions.
\end{grammarlens}

\subsection*{Your turn}
\begin{itemize}
\raggedright
  \item \emph{Say it:} \emph{kuṟukkiṭu}
  \item \emph{Say it:} \emph{kuṟukkiṭu}, once more
  \item \emph{Say it:} \emph{nāṉ pēcumpōtu avaṉ kuṟukkiṭṭāṉ.} — He interrupted while I was speaking.
  \item \emph{Recall:} say the Tamil for to talk over, then the Tamil for negotiations, then say \emph{kuṟukkiṭu} again
\end{itemize}

\begin{tcolorbox}[breakable,colback=teal!4,colframe=teal!35!black,title={Before you move on}]
What does \ta{குறுக்கிடு} mean? (\textquotedblleft{}To interrupt; to intervene\textquotedblright{}.) Say it once more.
\end{tcolorbox}
\section[\texorpdfstring{\ta{சமரசம்} (camaracam)}{சமரசம் (camaracam)}]{\ta{சமரசம்} (camaracam) — a compromise}

\label{lesson:TA-C356-camaracam}



\begin{quote}
Before the new one: say the Tamil for negotiations, then the Tamil for to interrupt.
\end{quote}

\subsection*{You'll want to know: \ta{சமரசம்}}
\textbf{\ta{சமரசம்}} — \emph{camaracam} — \textquotedblleft{}a compromise\textquotedblright{}.

\begin{quote}
\ta{நீண்ட} \ta{பேச்சுக்குப்} \ta{பிறகு} \ta{இருவரும்} \ta{சமரசத்துக்கு} \ta{வந்தார்கள்}.
\end{quote}

\begin{quote}
\emph{nīṇṭa pēccukkup piṟaku iruvarum camaracattukku vantārkaḷ.}
\end{quote}

\begin{quote}
After a long talk, the two reached a compromise.
\end{quote}

\ta{சமரசத்துக்கு} \ta{வா}, to come to a compromise; \ta{சமரசம்} \ta{செய்}, to make peace.

\begin{cousinweb}
\textbf{\ta{சமரசம்}} (\emph{camaracam}) is borrowed from Sanskrit \emph{samarasa}, \textquotedblleft{}having equal feelings\textquotedblright{} (Monier-Williams), from \emph{sama}, \textquotedblleft{}equal\textquotedblright{}, and \emph{rasa}, \textquotedblleft{}taste, feeling\textquotedblright{}: a state where both sides feel alike.
\end{cousinweb}

\begin{grammarlens}[title={What you've built}]
One more word for giving reasons and opinions.
\end{grammarlens}

\subsection*{Your turn}
\begin{itemize}
\raggedright
  \item \emph{Say it:} \emph{camaracam}
  \item \emph{Say it:} \emph{camaracam}, once more
  \item \emph{Say it:} \emph{nīṇṭa pēccukkup piṟaku iruvarum camaracattukku vantārkaḷ.} — After a long talk, the two reached a compromise.
  \item \emph{Recall:} say the Tamil for negotiations, then the Tamil for to interrupt, then say \emph{camaracam} again
\end{itemize}

\begin{tcolorbox}[breakable,colback=teal!4,colframe=teal!35!black,title={Before you move on}]
What does \ta{சமரசம்} mean? (\textquotedblleft{}A compromise\textquotedblright{}.) Say it once more.
\end{tcolorbox}
\section[\texorpdfstring{\ta{வாக்குவாதம்} (vākkuvātam)}{வாக்குவாதம் (vākkuvātam)}]{\ta{வாக்குவாதம்} (vākkuvātam) — a heated argument; a quarrel}

\label{lesson:TA-C356-vakkuvatam}



\begin{quote}
Before the new one: say the Tamil for to interrupt, then the Tamil for a compromise.
\end{quote}

\subsection*{You'll want to know: \ta{வாக்குவாதம்}}
\textbf{\ta{வாக்குவாதம்}} — \emph{vākkuvātam} — \textquotedblleft{}a heated argument; a quarrel\textquotedblright{}.

\begin{quote}
\ta{சின்ன} \ta{விஷயத்துக்காக} \ta{அவர்களிடையே} \ta{வாக்குவாதம்} \ta{ஏற்பட்டது}.
\end{quote}

\begin{quote}
\emph{ciṉṉa viṣayattukkāka avarkaḷiṭaiyē vākkuvātam ēṟpaṭṭatu.}
\end{quote}

\begin{quote}
An argument broke out between them over a small thing.
\end{quote}

\ta{வாக்குவாதம்} \ta{செய்}, to argue heatedly.

\begin{cousinweb}
\textbf{\ta{வாக்குவாதம்}} (\emph{vākkuvātam}) is built from two Sanskrit loans: \emph{vākku}, \textquotedblleft{}word\textquotedblright{}, from \emph{vāc}, \textquotedblleft{}speech\textquotedblright{} (Monier-Williams), and \emph{vātam}, \textquotedblleft{}argument\textquotedblright{}, from \emph{vāda} (CDIAL 11507): a war of words.
\end{cousinweb}

\begin{grammarlens}[title={What you've built}]
One more word for giving reasons and opinions.
\end{grammarlens}

\subsection*{Your turn}
\begin{itemize}
\raggedright
  \item \emph{Say it:} \emph{vākkuvātam}
  \item \emph{Say it:} \emph{vākkuvātam}, once more
  \item \emph{Say it:} \emph{ciṉṉa viṣayattukkāka avarkaḷiṭaiyē vākkuvātam ēṟpaṭṭatu.} — An argument broke out between them over a small thing.
  \item \emph{Recall:} say the Tamil for to interrupt, then the Tamil for a compromise, then say \emph{vākkuvātam} again
\end{itemize}

\begin{tcolorbox}[breakable,colback=teal!4,colframe=teal!35!black,title={Before you move on}]
What does \ta{வாக்குவாதம்} mean? (\textquotedblleft{}A heated argument; a quarrel\textquotedblright{}.) Say it once more.
\end{tcolorbox}
